\documentclass{exam}
\usepackage{../../shah311}

\hypersetup{colorlinks=true, linktoc=section, linkcolor=blue}

\pagestyle{headandfoot}
\firstpageheadrule
\runningheadrule
\firstpageheader{Prof. Rong \\ Real Analysis}{Homework 3}{Jeevan Shah}
\runningheader{Real Analysis \\ Homework 3}{}{Shah}
\firstpagefooter{}{}{}
\runningfooter{ }{\thepage}{ }

\printanswers
% hw3: pp47: 2.27, 2.3.2, 2.3.5, 2.3.6, 2.3.8, 2.3.11

% pp59: 2.4.1, 2.4.2, 2.4.3, 2.4.8 
\begin{document}
\colorbox{red}{\underline{$2.2.7$}:}
\begin{solution}
    \begin{parts}
        \part The sequence $(-1)^{n}$ is frequently in the set $\braces{1}$ since for every $N \in \NN$ there exists $n \geq N$ such that $a_n \in \braces{1}$, namely $n = N$ if $N$ is even or $n = N+1$ if $N$ is odd.

        \part Eventually is a stronger definition since a sequence that is eventually in a set is also frequently in a set. To see why, assume that a sequence $(a_n)$ is eventually in a set $A \subseteq \RR$. That is, there exists an $N_1 \in \NN$ such that $a_n \in A$ for all $n \geq N_1$. Then $(a_n)$ is frequently in $A$ since for every $N_2 \in \NN$ there exists $n \geq N_2$, namely $n \geq \max\braces{N_1, N_2}$, such that $a_n \in A$.
        
        \part A sequence $(a_n)$ converges to $a$ if $(a_n)$ is eventaully in the $\epsilon$-neighbourhood near $a$.

        \part If an infinite number of terms of a sequence $(x_n)$ are equal to $2$, then $(x_n)$ if frequently in the interval $(1.9, 2.1)$. It is not always eventually in the interval since an equal amount of terms maybe be equal to some $x \not\in (1.9, 2.1)$.
    \end{parts}
\end{solution}

\colorbox{red}{\underline{$2.3.2$}:}
\begin{solution}
    \begin{parts}
        \part \begin{proof}
            We must find $N \in \NN$ such that for all $n \geq N$, 
            \[
                \left|\frac{2x_n - 1}{3} - 1\right| < \epsilon.
            \]
            However, since $(x_n) \to 2$, for every $\epsilon > 0$, there exists $N_1 \in \NN$ such that $n_1 \geq N_1$ implies $|x_{n_1} - 2| < \frac{\epsilon}{2}$. Take $n \geq N \geq N_1$, then for any $\epsilon > 0$, we have 
            \begin{align*}
                \left|\frac{2x_n - 1}{3} - 1\right| &= \left|\frac{2x_n -1}{3} - \frac{3}{3}\right| \\
                &= \left|\frac{2x_n - 4}{3}\right| \\
                &= \frac{2}{3}\left|x_n - 2\right| < \frac{2}{3}\left(\frac{\epsilon}{2}\right) = \frac{\epsilon}{3} \\ 
                &\Rightarrow \left|\frac{2x_n - 1}{3} - 1\right| < \frac{\epsilon}{3} < \epsilon.
            \end{align*}
            Thus, the sequence converges.
        \end{proof}

        \part \begin{proof}
            We must show that for any $\epsilon > 0$ there exists $N \in \NN$ such that $n \geq N$ implies 
            \[
                \left|\frac{1}{x_n} - \frac{1}{2}\right| < \epsilon.
            \]
            However, since $(x_n) \to 2$, there exists $N_1 \in \NN$ such that $n \geq N_1$ implies that $|x_n - 2| < \epsilon$. There must also exist $N_2 \in \NN$ such that $n \geq N_2$ implies that
            \[
                |x_n - 2| < 1 \Rightarrow 1 < |x_n| < 3 \Rightarrow 1 < |x_n|. 
            \]
            Choose $N = \max\braces{N_1, N_2}$ so for $n \geq N$ 
            \[
                \left|\frac{1}{x_n} - \frac{1}{2}\right| = \left|\frac{2-x_n}{2x_n}\right| < \frac{|x_n - 2|}{2} < \frac{\epsilon}{2} < \epsilon.
            \]
            It follows that the sequence converges. 
        \end{proof}
    \end{parts}
\end{solution}

\newpage

\colorbox{red}{\underline{$2.3.5$}:}
\begin{solution}
    \begin{proof}
        $(\Rightarrow)$ Assume that $(x_n)$ and $(y_n)$ are both convergent with $\lim x_n = \lim y_n = L$. Then there exists $N_1, N_2 \in \NN$ such that for all $n_1 \geq N_1$ and $n_2 \geq N_2$ 
        \[
            |x_{n_1} - L| < \epsilon \quad\text{and}\quad |y_{n_2} - L| < \epsilon.
        \]
        Take $N = \max{N_1, N_2}$. Then $(z_n) \to L$ since 
        \[
            |z_n - L| < \epsilon 
        \]
        for all $\epsilon > 0$ because regardless if $z_n = x_n$ or $z_n = y_n$, the inequality holds true. \\
        $(\Leftarrow)$ Assume that $(z_n) \to L$. That is, for every $\epsilon > 0$, there exists $N \in \NN$ such that $n \geq N$ implies $|z_n - L| < \epsilon$. However, since $z_n$ is either $x_n$ or $y_n$ we know that for $n \geq N$
        \[
            |x_{n} - L| < \epsilon \quad\text{and}\quad |y_{n} - L| < \epsilon.
        \]
        Thus, by definition we must have that $(x_n) \to L$ and $(y_n) \to L$.
    \end{proof}
\end{solution}

\colorbox{red}{\underline{$2.3.6$}:}
\begin{solution}
    \begin{proof}
        We start by multiplying by the conjugate to see 
        \begin{align*}
            n - \sqrt{n^2 + 2n}\left(\frac{n + \sqrt{n^2 + 2n}}{n + \sqrt{n^2 + 2n}}\right) &= \frac{n^2 - (n^2 + 2n)}{n + \sqrt{n^2 + 2n}} \\
            &= \frac{2n}{n + \sqrt{n^2\left(1 + \frac{2}{n}\right)}} \\
            &= \frac{2n}{n + n\sqrt{1 + \frac{2}{n}}} \\
            &= \frac{2}{1 + \sqrt{1 + \frac{2}{n}}}
        \end{align*}
        Using the given fact that $\lim (\frac{1}{n}) = 0$ and that if $(x_n) \to 0$ then $(\sqrt{x_n}) \to 0$, 
        \[
            \lim (n - \sqrt{n^2 + 2n}) = \lim\left(\frac{2}{1 + \sqrt{1 + \frac{2}{n}}}\right) = \frac{2}{1 + 1} = 1.
        \]
        Thus, the limit exists and is equal to $1$.
    \end{proof}
\end{solution}

\newpage 

\colorbox{red}{\underline{$2.3.8$}:}
\begin{solution}
    \begin{parts}
        \part \begin{proof} Let $(x_n) \to x$ and $p(x)$ be a polynomial. Any polynomial $p(x)$ of degree $n$ can be represented in the form 
        \[
            p(x) = a_{n}x^{n} + a_{n-1}x^{n-1} + \cdots + a_{1}x + a_{0}. 
        \] 
        Then, 
        \begin{align*}
            \lim p(x_n) &= \lim (a_{n}x_n^{n} + a_{n-1}x_n^{n-1} + \cdots + a_{1}x_n + a_{0}) \\
            &= \lim(a_n x_n^n) + \lim(a_{n-1}x_n^{n-1}) + \cdots + \lim(a_1 x_n) + \lim a_0 \\
            &= a_n \lim x_n^n + a_{n-1}\lim x_n^{n-1} + \cdots + a_1\lim x_n + a_0.
        \end{align*}
        But, 
        \[
            \lim x_n^n = \underbrace{\lim x_n \cdot \lim x_n \cdots \lim x_n}_{n \text{ times}}.
        \]
        Since $(x_n) \to x$ it is equivalent to say that $\lim x_n = x$, so then $\lim x_n^{n} = x^{n}$ by the algebraic limit theorem. It follows that 
        \[
            \lim p(x_n) = a_n x^{n} + a_{n-1}x^{n-1} + \cdots + a_{1}x + a_0 = p(x).
        \]
        Then $p(x_n) \to p(x)$.
    \end{proof}

    \part Consider the sequence $x_n = \frac{1}{n}$ and the function 
    \[
        f(x) = \begin{cases}
            1 &\quad \text{ if } x \neq 0 \\
            0 &\quad \text{ if } x = 0
        \end{cases}.
    \]
    Then, the sequence $f(x_n)$ converges since it is a constant sequence of only $1$. But $x_n \to 0$ and $f(0) = 0$ so $f(x_n)$ converges to a value not equal to $f(x)$.
    \end{parts}
\end{solution}

\colorbox{red}{\underline{$2.3.11$}:}
\begin{solution}
\begin{parts}
    \part \begin{proof}
        We must show that if $(x_n) \to L$ then $(y_n) \to L$. That if, for any $\epsilon > 0$ there must exist $N \in \NN$ such that $n \geq N$ implies $|y_n - L| < \epsilon$. From $(x_n) \to L$ we know that there exists $N_1 \in \NN$ such that $n \geq N_1 \Rightarrow |x_n - N_1| < \epsilon$. Consider
        \begin{align*}
            y_n - L &= \frac{1}{n}\sum_{k=1}^{n}(x_n - L) \\
            &= \frac{1}{n}\left(\sum_{k=1}^{N_1 - 1} (x_k - L) + \sum_{k = N_1}^{n} (x_k - L)\right) \\
            &\leq \frac{1}{n}\left(\left|\sum_{k=1}^{N_1 - 1}(x_k - 1)\right| + \sum_{k=N_1}^{n}|x_k - L|\right).
        \end{align*}
        However, as $n \to \infty$ then $\frac{1}{n} \to 0$, so for $n$ large 
        \[
            \frac{1}{n}\left|\sum_{k=1}^{N_1 - 1} (x_k - 1)\right| \to 0.
        \]
        As well, since $(x_n) \to L$ then $|x_k - L| < \epsilon$. So, 
        \[
            \sum_{k=N_1}^{n} |x_k - L| < (n - N_1 +1)\epsilon \Rightarrow \frac{1}{n}\sum_{k=N_1}^{n} |x_k - L| < \frac{n - N_1 + 1}{n}\epsilon
        \]
        since there are $n - N_1 + 1$ terms being summed. For $n$ sufficently large we have 
        \[
            |y_n - L| \leq \frac{1}{n}\left|\sum_{k=1}^{N_1 - 1} x_k - L\right| + \frac{1}{n}\sum_{k=N_1}^{n}|x_k - L| \leq \frac{n - N_1 + 1}{n}\epsilon < \epsilon.
        \]
        So, $(y_n) \to L$.
    \end{proof}

    \part Consider $x_n = (-1)^n$. Then $(x_n)$ does not converge but 
    \[
        y_n = \frac{x_1 + x_2 + \cdots + x_n}{n} 
    \]
    will converge to $0$ since the numerate will alternate between $0$ (for even values of $n$) and $-1$ (for odd values of $n$). Thus, as $n \to \infty$ the denominator also tends to $\infty$ so the average tends to $0$. Thus, 
    \[
        \lim y_n = 0 
    \]
    and so $y_n$ is convergent.
\end{parts}
\end{solution}

\colorbox{red}{\underline{$2.4.1$}:}
\begin{solution}
    \begin{parts}
        \part We will show that $(x_n)$ is bounded and monotone and thus convergent. \\
        \underline{Bounded:} We will show $x_n \leq 3$ for all $n$ using induction. The base case is trivial as $x_1 = 3 \leq 3$. Now we assume that $x_n \leq 3$ and we must show that $x_{n+1} \leq 3$. But 
        \[
            x_{n+1} = \frac{1}{4 - x_n} \leq \frac{1}{4 - 3} = 1 \leq 3 
        \]
        which was to be shown. \\
        \underline{Monotone:} We will show that $x_n \geq x_{n+1}$ for all $n$ by induction. The base case holds true since $x_1 = 3 \geq 1 = x_2$. Now, we assume that $x_{n+1} \geq x_{n}$ and we must show that $x_{n} \geq x_{n+1}$. But 
        \[
            x_{n+1} = \frac{1}{4 - x_n} \leq \frac{1}{4-x_{n-1}} = x_{n}
        \]
        which was to be shown.

        \part Since $\lim x_n$ exists, we may say that $\lim x_n = x$, or equivalently, $(x_n) \to x$. Then for all $\epsilon > 0$ there exists $N \in \NN$ such that $n \geq N \Rightarrow |x_n - x| < \epsilon$. Thus, since $n + 1 > n, |x_{n+1} - x| < \epsilon$ so $\lim x_{n+1} = \lim x_n = x$.

        \part We start by taking the limit of both sides and calculating 
        \begin{align*}
            x_{n+1} = \frac{1}{4 - x_n} &\to x = \frac{1}{4 - x} \\
            &\Rightarrow x(4 - x) = 1 \\
            &\Rightarrow x^2 - 4x + 1 = 0 \\
            &\Rightarrow x = \frac{4 \pm \sqrt{16 - 4}}{2} \\
            &\Rightarrow x = 2 \pm \sqrt{3}.
        \end{align*}
        Since $(x_n)$ is bounded by $3$ and $2 + \sqrt{3} > 3$, 
        \[
            \lim x_n = 2 - \sqrt{3}.
        \]
    \end{parts}
\end{solution}

\colorbox{red}{\underline{$2.4.2$}:}
\begin{solution}
    \begin{parts} 
        \part This argument assumes the existance of the limit which is false. The sequence will alternate between $1$ and $2$ forever and thus the limit does not exist.

        \part We may apply the strategy used in $(a)$ if $y_n$ is bounded and monotone. \\ 
        \underline{Bounded:} We claim that $|y_n| \leq 3$ for all $n$. We will show this by induction. The base case is trivial as $y_1 = 1 \leq 3$. Now, we assume $y_n \leq 3$ and we must show that $y_{n+1} \leq 3$. But 
        \[
            y_{n+1} = 3 - \frac{1}{y_n} \leq 3 - \frac{1}{3} < 3 
        \] 
        which was to be shown. \\
        \underline{Monotone:} We claim that $y_n \leq y_{n+1}$ for all $n$. We will show this by induction. The base is that $y_1 = 1 \leq 2 = y_2$ which is true. Now, we assume that $y_{n-1} \leq y_n$ and we must show that $y_n \leq y_{n+1}$. But, 
        \[
            y_{n+1} = 3 - \frac{1}{y_n} \geq 3 - \frac{1}{y_{n-1}} = y_n 
        \]
        which was to be shown. \\
        Thus, since $y_n$ is bounded and monotone, the limit exists.
    \end{parts}
\end{solution}

\colorbox{red}{\underline{$2.4.3$}:}
\begin{solution}
    \begin{parts}
        \part We will show that $a_{n+1} = \sqrt{2 + a_n}, a_1 = \sqrt{2}$ is bounded and montone and thus convergent. \\ 
        \underline{Bounded:} We claim that $|a_n| \leq 2$ for all $n$. We will show this using induction. The base case is trivial as $a_1 = \sqrt{2} < 2$. Now, we assume that $a_n < 2$ and we must show that $a_{n+1} < 2$. But 
        \[
            a_{n+1} = \sqrt{2 + a_n} < \sqrt{2 + 2} = 2 
        \]
        which was to be shown. \\
        \underline{Monotone:} We will show that $a_n < a_{n+1}$ for all $n$ using induction. The base case is true since $a_1 = \sqrt{2} < \sqrt{2 + \sqrt{2}} = a_1$. Now, we assume that $a_{n-1} < a_n$ and we must show that $a_n < a_{n+1}$. But 
        \[
            a_{n+1} = \sqrt{2 + a_n} \geq \sqrt{2 + a_{n-1}} = a_n 
        \]
        which was to be shown. \\
        Now since $\lim a_{n+1} = \lim a_n$ we can find that 
        \begin{align*}
            a_{n+1} = \sqrt{2 + a_n} \to a = \sqrt{2 + a} &\Rightarrow a^2 = 2 + a \\
            &\Rightarrow a^2 - a - 2 = 0 \\
            &\Rightarrow (a+1)(a-2) = 0 \\
            &\Rightarrow a = 2, -1.
        \end{align*}
        Since $a_n \geq 0$ for all $n$, $a \not< 0$ so $a = 2$ and $(a_n) \to 2$.

        \part We will show that $a_{n+1} = \sqrt{2a_n}, a_1 = \sqrt{2}$ is bounded and monotone and thus convergent. \\
        \underline{Bounded:} We claim that $a_n \leq 2$ for all $n$. We will show this using induction. The base case is trivial as $a_1 = \sqrt{2} \leq 2$. Now, we assume that $a_n \leq 2$ and we must show that $a_{n+1} \leq 2$. But, 
        \[
            a_{n+1} = \sqrt{2 a_n} \leq \sqrt{2 \cdot 2} = 2
        \]
        which was to be shown. \\
        \underline{Monotone:} We will show that $a_{n+1} \geq a_n$ using induction. The base case is true since $a_1 = \sqrt{2} \leq \sqrt{2\sqrt{2}} = a_2$. Now, we assume that $a_{n-1} \leq a_n$ and we must show that $a_{n} \leq a_{n+1}$. But, 
        \[
            a_{n+1} = \sqrt{2 a_n} \geq \sqrt{2a_{n-1}} = a_n 
        \]
        which was to be shown. \\
        Thus, since the limit exists and $\lim a_{n+1} = \lim a_n$ we can find that 
        \begin{align*}
            a_{n+1} = \sqrt{2 a_n} \to a = \sqrt{2a} &\Rightarrow a^2 = 2a \\
            &\Rightarrow a^2 - 2a = 0 \\
            &\Rightarrow a(a - 2) = 0 \\
            &\Rightarrow a = 0, 2.
        \end{align*}
        Since $a_n \geq 0$ for all $n$, $a \neq 0$ and thus $a = 2$ so $(a_n) \to 2$.
    \end{parts}
\end{solution}

\colorbox{red}{\underline{$2.4.8$}:}
\begin{solution}
    \begin{parts}
        \part Considering the first three partial sums we can see that 
        \begin{align*}
            S_1 &= \frac{1}{2} \\
            S_2 &= \frac{1}{2} + \frac{1}{4} = \frac{3}{4} \\
            S_3 &= \frac{3}{4} + \frac{7}{8}.
        \end{align*}
        We may now conclude that 
        \[
            S_n = 1 - \frac{1}{2^n}.
        \]
        We can now show that $S_n \leq 1$ for all $n$. Since for any $n$,
        \[
            \frac{1}{2^n} > 0 
        \]
        it follows that 
        \[
            1 - \frac{1}{2^n} \leq 1 
        \]
        which was to be shown. We will now show that $S_n$ is monotone, that is, for all $n, S_n < S_{n+1}$. Since $n \in \NN$ we may conclude that $n < n+1$ for all $n$ so, 
        \[
            \frac{1}{2^{n}} \geq \frac{1}{2^{n+1}}.
        \]
        It follows that 
        \[
            S_n = 1 - \frac{1}{2^n} \leq 1 - \frac{1}{2^{n+1}} = S_{n+1}.
        \]
        Since $S_n$ is bounded and montone, it follows that it is convergent. 
        \part Considering the first three partial sums we can see that 
        \begin{align*}
            S_1 &= \frac{1}{2} \\
            S_2 &= \frac{1}{2} + \frac{1}{6} = \frac{2}{3} \\
            S_3 &= \frac{2}{3} + \frac{1}{12} = \frac{3}{4}.
        \end{align*}
        We may now conclude that 
        \[
            S_n = \frac{n}{n+1}.
        \]
        We can now show that $S_n \leq 1$ for all $n$. Since $n \in \NN$ it is always true that $n < n+1$. Thus 
        \[
            \frac{n}{n+1} < 1 
        \]
        for all $n$. We will now show that $S_n$ is monotone, that is, for all $n$ we have that $S_n < S_{n+1}$. But 
        \[
            S_n = \frac{n}{n+1} < \frac{n+1}{n+2} = S_{n+1}
        \]
        since 
        \begin{align*}
            \frac{n}{n+1} < \frac{n+1}{n+2} &\Rightarrow n(n+2) < (n+1)^2 \\
            &\Rightarrow n^2 + 2n < n^2 + 2n + 1 \\
            &\Rightarrow 0 < 1
        \end{align*}
        which is true. Thus $S_n < S_{n+1}$. Since $S_n$ is bounded and montone, it must also be convergent.

        \part Considering the first three partial sums we can see that 
        \begin{align*}
            S_1 &= \log(2) \\
            S_2 &= \log(2) + \log\left(\frac{3}{2}\right) = \log(3) \\
            S_3 &= \log(3) + \log\left(\frac{4}{3}\right) = \log(4).
        \end{align*}
        We may now conclude that 
        \[
            S_n = \log(n+1).
        \]
        It is clear that $S_n$ is an unbounded sequence and so it must diverge.
    \end{parts}
\end{solution}
\end{document}